\documentclass[10pt]{article}
\usepackage[T1]{fontenc}
\usepackage[utf8]{inputenc}
\usepackage[french]{babel}
\usepackage{amsmath,amssymb,mathrsfs,mathpazo}
\usepackage{geometry}
\geometry{a4paper,margin=2cm}
\usepackage{graphicx,array,tabularx,tabularray,colortbl,multirow,multicol,booktabs,enumitem,pifont,ulem,fancybox,diagbox,cancel}
\usepackage{pstricks,pstricks-add,pst-plot,pst-tree,pst-node,pst-text,pst-func,pst-math,pst-bspline,pst-3dplot,pst-eucl,pst-solides3d}
\let\pscalculate\undefined
\usepackage{pst-bezier}
\usepackage{tikz,pgfplots,tkz-tab}
\usetikzlibrary{arrows,arrows.meta,patterns,calc,positioning,trees,decorations.markings,shapes.geometric,shapes.misc,shapes.arrows,intersections,decorations.pathmorphing,decorations.pathreplacing,backgrounds,shadows,babel}
\pgfplotsset{compat=1.18}
\usepgfplotslibrary{fillbetween}
\usepackage[np]{numprint}
\usepackage{listings,makecell,pas-tableur,twoopt}
\let\barre\undefined
\usepackage[pstricks]{tablvar}
\definecolor{darkgreen}{RGB}{0,100,0}
\providecommand{\red}{\color{red}}
\providecommand{\blue}{\color{blue}}
\providecommand{\green}{\color{green}}
\providecommand{\black}{\color{black}}
\providecommand{\white}{\color{white}}
\providecommand{\gray}{\color{gray}}
\providecommand{\cyan}{\color{cyan}}
\providecommand{\magenta}{\color{magenta}}
\providecommand{\yellow}{\color{yellow}}
\providecommand{\violet}{\color{violet}}
\providecommand{\orange}{\color{orange}}
\providecommand{\euro}{\texteuro}
\providecommand{\EUR}{\texteuro}
\providecommand{\ieme}{\textsuperscript{e}}
\providecommand{\degres}{\ensuremath{^\circ}}
\providecommand{\degree}{\ensuremath{^\circ}}
\providecommand{\vv}[1]{\overrightarrow{#1}}
\providecommand{\Coord}[1]{#1}
\providecommand{\decosix}{}
\providecommand{\starredbullet}{\ensuremath{\bullet}}
\pagestyle{empty}
\setlength{\parindent}{0pt}
\usepackage{adjustbox}\DeclareUnicodeCharacter{2236}{:}\DeclareUnicodeCharacter{2190}{\ensuremath{\leftarrow}}
\providecommand{\R}{\mathbb{R}}
\providecommand{\N}{\mathbb{N}}
\providecommand{\Z}{\mathbb{Z}}
\providecommand{\Q}{\mathbb{Q}}
\providecommand{\D}{\mathbb{D}}
\providecommand{\C}{\mathbb{C}}
\providecommand{\e}{{\mathrm{e}}}
\providecommand{\vect}[1]{\overrightarrow{#1}}
\providecommand{\vectt}[1]{\overrightarrow{\mathrm{#1}}}
\providecommand{\vvt}[1]{\overrightarrow{\mathrm{#1}}}
\providecommand{\np}[2][]{{#2~#1}}
\providecommand{\vv}[1]{{\overrightarrow{#1}}}
\renewcommand{\d}{{\mathrm{d}}}
\providecommand{\degres}{{^\circ}}
\providecommand{\degre}{{^\circ}}
\providecommand{\blue}{\color{blue}}
\providecommand{\red}{\color{red}}
\providecommand{\enskip}{\quad}
\providecommand{\llbracket}{\lbrack\!\lbrack}
\providecommand{\rrbracket}{\rbrack\!\rbrack}
\providecommand{\euro}{{\text{€}}}
\providecommand{\no}{{\text{n°}}}
\providecommand{\interval}[2]{{[#1;#2]}}
\providecommand{\intervFO}[2]{{[#1;#2[}}
\providecommand{\dotfill}{\ldots}

\providecommand{\ds}{\displaystyle}
\providecommand{\textding}[1]{#1}
\providecommand{\Oij}{\(\mathrm{O};\vec{\imath},\vec{\jmath}\)}
\providecommand{\Oijk}{\(\mathrm{O};\vec{\imath},\vec{\jmath},\vec{k}\)}
\providecommand{\Ouv}{\(\mathrm{O};\vec{u},\vec{v}\)}
\providecommand{\intervFF}[2]{\ensuremath{[#1;#2]}}
\renewcommand{\euro}{\ifmmode\text{\texteuro}\else\texteuro\fi}
\DeclareUnicodeCharacter{00B0}{\ifmmode^{\circ}\else\textdegree\fi}
\DeclareUnicodeCharacter{20AC}{\ifmmode\text{\texteuro}\else\texteuro\fi}
% Énoncés : Éducation nationale. Transcription et corrigés : APMEP, Denis Vergès et contributeurs. Découpage et classement : CoopMaths, Sébastien Lozano, Éric Elter et contributeurs. Adaptation de la présentation : Maths Explorer.
% CC BY-SA : https://coopmaths.fr/www/about/
% https://coopmaths.fr/alea/static/bac/2021/tex/bac_2021_03_sujet2_5.tex
\begin{document}
\noindent\textit{APMEP / CoopMaths. CC BY-SA. Présentation adaptée par Maths Explorer.}\par\medskip
\medskip

\begin{tabularx}{\linewidth}{|X|}\hline
\textbf{Principaux domaines abordés :
Fonction logarithme ; dérivation.}\\ \hline
\end{tabularx}

\bigskip

\textbf{Partie I: Étude d'une fonction auxiliaire}

\medskip

Soit $g$ la fonction définie sur $]0~;~+\infty[$ par :

\[g(x) = \ln(x) + 2x - 2.\]

\smallskip

\begin{enumerate}
\item Déterminer les limites de $g$ en $+\infty$ et $0$.
\item Déterminer le sens de variation de la fonction $g$ sur $]0~;~ +\infty[$.
\item Démontrer que l'équation $g(x) = 0$ admet une unique solution $\alpha$ sur $]0~;~ +\infty[$. 
\item Calculer $g(1)$ puis déterminer le signe de $g$ sur $]0~;~ +\infty[$.
\end{enumerate}

\bigskip

\textbf{Partie II : Étude d'une fonction } \boldmath $f$\unboldmath

\medskip

On considère la fonction $f$, définie sur $]0~;~ +\infty[$ par: 

\[f(x) = \left(2 - \dfrac{1}{x}\right)\left(\ln (x) - 1\right).\]

\begin{enumerate}
\item 
	\begin{enumerate}
		\item On admet que la fonction $f$ est dérivable sur $]0~;~ +\infty[$ et on note $f'$ sa dérivée.

Démontrer que, pour tout $x$ de $]0~;~ +\infty[$, on a : 

\[f'(x) = \dfrac{g(x)}{x^2}.\]

		\item Dresser le tableau de variation de la fonction $f$ sur $]0~;~ +\infty[$. Le calcul des limites n'est pas demandé.
	\end{enumerate}
\item  Résoudre l'équation $f(x) = 0$ sur $]0~;~ +\infty[$ puis dresser le tableau de signes de $f$ sur l'intervalle $]0~;~ +\infty[$.
\end{enumerate}

\bigskip

\textbf{Partie III : Étude d'une fonction \boldmath $F$\unboldmath{} admettant pour dérivée la fonction \boldmath $f$\unboldmath}

\medskip

On admet qu'il existe une fonction $F$ dérivable sur $]0~;~ +\infty[$ dont la dérivée $F'$ est la fonction $f$.

Ainsi, on a : $F' = f$.

On note $\mathcal{C}_F$ la courbe représentative de la fonction $F$ dans un repère orthonormé \Oij. On ne cherchera pas à déterminer une expression de $F(x)$.

\medskip

\begin{enumerate}
\item Étudier les variations de $F$ sur $]0~;~ +\infty[$.
\item La courbe $\mathcal{C}_F$ représentative  de $F$ admet-elle des tangentes parallèles à l'axe des abscisses ?

Justifier la réponse.
\end{enumerate}

\end{document}
